\documentclass[11pt]{article}
\usepackage[a4paper,margin=1in]{geometry}
\usepackage{amsmath,amssymb,mathtools,bm}
\usepackage{enumitem,booktabs,array}
\usepackage{tcolorbox}
\usepackage{hyperref}
\usepackage[protrusion=true,expansion=false]{microtype}
\hypersetup{colorlinks=true,linkcolor=blue,urlcolor=blue,citecolor=blue}
\setlist[itemize]{topsep=2pt,itemsep=2pt}
\setlist[enumerate]{topsep=2pt,itemsep=2pt}
\newtcolorbox{keybox}{colback=blue!3!white,colframe=blue!40!black,boxrule=0.4pt,arc=1mm}
\newtcolorbox{alertbox}{colback=red!3!white,colframe=red!40!black,boxrule=0.4pt,arc=1mm}
\newtcolorbox{answerbox}{colback=green!3!white,colframe=green!40!black,boxrule=0.4pt,arc=1mm}
\newtcolorbox{drillbox}{colback=yellow!5!white,colframe=orange!60!black,boxrule=0.4pt,arc=1mm}
\newcommand{\EV}{\mathbb{E}}
\newcommand{\Var}{\mathrm{Var}}
\newcommand{\Cov}{\mathrm{Cov}}
\newcommand{\blank}[1]{\underline{\hspace{#1}}}

\title{W08-03: Acharya, Lochstoer \& Ramadorai (2013, JFE) \\
\large ``Limits to Arbitrage and Hedging: Evidence from Commodity Markets'' --- Active Recall Workbook}
\author{Banking \& Risk Management --- Exam Prep}
\date{\today}
\begin{document}\maketitle

\begin{keybox}
\textbf{Paper in one sentence.} When commodity producers face capital constraints (measured by their \emph{default probability}), they hedge less; futures markets then clear at lower prices (higher risk premia) --- empirically confirming a \emph{producer-hedging pressure} theory of commodity risk premia. Data span 1990--2008 across crude oil, gas, and other commodities.

\textbf{Placement.} Connects the risk-management-motives literature (Froot--Scharfstein--Stein, AW, JJ, GT) to the asset-pricing literature on commodities (Hirshleifer, Hong--Yogo, Gorton--Rouwenhorst).
\end{keybox}

\section*{Round 1: Complete Walk-Through}

\subsection*{1.1 Theory}
If producers want to hedge their future output but speculators require risk premium to take the other side, futures prices are below expected spot (normal backwardation). When producers are capital-constrained, they hedge \emph{less}, net speculator demand \emph{falls}, futures-spot basis \emph{widens}, and realised returns to speculators \emph{rise} --- i.e., the hedging premium is endogenous in producer distress.

\subsection*{1.2 Empirical design}
Measure producer capital constraints with default-probability proxy (Moody's KMV $\text{EDF}$) at a commodity-specific level (aggregate across producers in oil, gas, etc.). Test:
\[
R^{\text{futures}}_{c,t+1}=\alpha_c + \beta\cdot \text{DefaultProb}_{c,t}+X_{c,t}'\gamma+\varepsilon_{c,t+1}.
\]
Expected $\beta>0$: higher producer default probability now predicts higher subsequent futures returns for commodity $c$.

\subsection*{1.3 Data}
Monthly futures returns on crude oil, natural gas, heating oil, and industrial metals, 1990--2008. Producer financials from Compustat. Producer default probabilities from Moody's KMV.

\subsection*{1.4 Headline results}
\begin{itemize}
\item Producer default probability predicts futures returns with the expected positive sign and statistically significant magnitude ($\sim 5$--$10\%$ annualised risk premium differential between high- and low-constraint periods).
\item Hedging activity (CFTC positions) falls when producers are more constrained.
\item Effect is strongest in energy commodities (crude, natural gas) where producers are concentrated.
\end{itemize}

\subsection*{1.5 Interpretation}
\begin{itemize}
\item Supports a \emph{limits-to-arbitrage} view of commodity risk premia: speculators must be compensated when hedging supply is weak.
\item Links corporate risk management directly to asset-pricing outcomes.
\item Distinct from a pure-storage or pure-speculative view of commodity markets.
\end{itemize}

\subsection*{1.6 Caveats}
\begin{alertbox}
\begin{itemize}
\item Default probability is proxied; noise attenuates estimates.
\item CFTC-based hedging-activity data aggregates across reporting categories with possible misclassification.
\item Endogeneity of default probability to commodity prices themselves (risk/return simultaneity); ALR use lagged measures and robustness checks.
\end{itemize}
\end{alertbox}

\section*{Round 2: Level 1 Fill-in}
\begin{enumerate}
\item Producer-constraint measure: \blank{2cm} default probability.
\item Expected sign on $\beta$: \blank{1cm}.
\item Hedging-activity data: \blank{1cm} positions.
\item Strongest effect sector: \blank{1.5cm}.
\item Sample period: \blank{1cm}--\blank{1cm}.
\end{enumerate}

\section*{Round 3: Level 2 Prompts}
\begin{enumerate}
\item State the producer-hedging-pressure story in words and equations.
\item What role does limits-to-arbitrage play?
\item Relate ALR to FSS (1993) and Allayannis--Weston (2001).
\item Propose a test of causality rather than predictability.
\end{enumerate}

\section*{Round 4: Minimal Scaffolding}
From a blank page: (i) theory, (ii) empirical spec, (iii) headline sign and magnitude, (iv) interpretation, (v) two caveats.

\section*{Round 5: Blank Page}
Essay: corporate hedging frictions and commodity risk premia.

\vspace{6pt}
\begin{drillbox}
\textbf{Speed Drill.} (1) Producer constraint measure. (2) Predicted sign. (3) Data category for hedging positions. (4) Asset-pricing implication. (5) Sample years.
\end{drillbox}

\section*{Quick Reference \& Answer Key}
\begin{answerbox}
\begin{itemize}
\item \textbf{Theory.} Constrained producers hedge less $\Rightarrow$ futures risk premium rises.
\item \textbf{Proxy.} Producer KMV default probability, commodity-specific.
\item \textbf{Result.} Positive predictive coefficient on futures returns; $5$--$10\%$ differential.
\item \textbf{Implication.} Corporate financial frictions show up in commodity risk premia.
\end{itemize}
\end{answerbox}

\begin{keybox}
\textbf{30-second exam pitch.} Acharya, Lochstoer \& Ramadorai (2013) bridge corporate risk management and asset pricing. The idea: when commodity producers are capital-constrained, they hedge less; with fewer natural sellers of futures, speculator risk premia must widen to clear the market. Using Moody's KMV default probabilities at the commodity-producer level as a constraint proxy, and CFTC-based hedging-position data, they show that producer distress positively predicts subsequent commodity futures returns, with magnitudes of 5--10\% annualised between high- and low-constraint periods. Effects are strongest in energy commodities where producer concentration is highest. The paper makes the limits-to-arbitrage / hedging-pressure theory of commodity risk premia operational and connects corporate-finance frictions to macro-asset-pricing outcomes.
\end{keybox}

\end{document}
